%% ma: L12-B14-0147
%% lop: 12
%% chuong: 5
%% bai: 14
%% dang: 12.14.44
%% loai: TLN
%% muc: VD
%% dap_an: "16"
%% nguon: "Ảnh giáo viên gửi 25/11/2026 (tài liệu ôn tập Phương trình mặt phẳng), Phần III Câu 4 (bốn điểm A(2;1;3), C(3;1;3), D(4;3;1), B(1;-2;1))"
%% kien_thuc: "Bài 14 (tích có hướng, mặt phẳng song song với đường thẳng)"
%% trang_thai: cho-duyet
%% ly_do: "Dạng toán mới đề xuất, chờ giáo viên duyệt"
%% ghi_chu: "Đáp án tự giải (sympy): 10x-4y+z-19=0."
\begin{ex}
Trong không gian $Oxyz$, cho bốn điểm $A(2;1;3)$, $C(3;1;3)$, $D(4;3;1)$, $B(1;-2;1)$. Mặt phẳng đi qua hai điểm $A,B$ và song song với đường thẳng $CD$ có phương trình dạng $10x+by+cz+d=0$. Tính $b+c-d$.
\shortans{16}
\loigiai{$\overrightarrow{AB}=(-1;-3;-2)$, $\overrightarrow{CD}=(1;2;-2)$, $\left[\overrightarrow{AB},\overrightarrow{CD}\right]=(10;-4;1)$ là một vectơ pháp tuyến. Mặt phẳng: $10(x-2)-4(y-1)+(z-3)=0\Leftrightarrow10x-4y+z-19=0$ (điểm $B$ thỏa mãn, $C\notin$ mặt phẳng). Vậy $b=-4$, $c=1$, $d=-19$ và $b+c-d=16$.}
\end{ex}
